\section[Absence as presence]{Absence as presence: the resurrection and exaltation of Christ}

\begin{quotex}
When the Divine Light penetrates the soul, it is united with God as light with light. This is the light of faith. Faith bears the soul to heights unreachable by her natural senses and faculties.

\flright{\textsc{Meister Eckhart}}
\end{quotex}

Resurrection and faith are a critical matter for Christians; especially in this time of Easter. More so, for traditionalists, it leads to a more profound understanding of the resurrection of the spirit through initiatic death and rebirth (i.e. the cross bearing, death and rebirth of Christ). In the Gospel of Luke, there are plenty of important facts that points out to something more than just merely a biography of an event, marked in time and space: it is a timeless and profound meaning of rebirth; as it is said in the liturgy of the Easter Vigil : it is the rebirth of the whole of creation.

In the New Testament, the writers have used words like «eigeirô» (to wake up) in response to the term «koïmao» (to sleep); «anistème» (to get up) and «anastasis» (to get up from a fall); and saying that he now lives («zaô»). Of course, those terms can be interpreted in a Gnostic mindset and deliver interesting information (getting up from a fall, waking up, etc.), but what is most important is the fact that they are not alone. If the author would only have used words like «living», it would have been only a matter of the resurrection of the flesh, a mere reanimation of a corpse. But that is not the case.

For Luke and the evangelists, we cannot reduce the resurrection to a matter of «living man». It is about a living God. That is why it was added that He is glorified and exalted, that He has ascended. The vocabulary chosen is always that of glorification (as with God) and verticality (getting up literally meaning to put oneself in a straight position). It is to enter the glory of God.

Now as for the story of Luke, there is hidden beauty in it, waiting to be unveiled both by examination and meditation; especially as this is the text used in Catholic Easter liturgy on Sunday. Let us remind us of the Douay-Rheims Bible\footnote{\url{https://www.drbo.org/chapter/49024.htm}}:

\begin{quotationx}
And behold, two of them went, the same day, to a town which was sixty furlongs from Jerusalem, named Emmaus. And they talked together of all these things which had happened. And it came to pass, that while they talked and reasoned with themselves, Jesus himself also drawing near, went with them.

But their eyes were held, that they should not know him. And he said to them: What are these discourses that you hold one with another as you walk, and are sad? And the one of them, whose name was Cleophas, answering, said to him: Art thou only a stranger to Jerusalem, and hast not known the things that have been done there in these days? To whom he said: What things? And they said: Concerning Jesus of Nazareth, who was a prophet, mighty in work and word before God and all the people; And how our chief priests and princes delivered him to be condemned to death, and crucified him.

But we hoped, that it was he that should have redeemed Israel: and now besides all this, today is the third day since these things were done. Yea and certain women also of our company affrighted us, who before it was light, were at the sepulchre, And not finding his body, came, saying, that they had also seen a vision of angels, who say that he is alive. And some of our people went to the sepulchre, and found it so as the women had said, but him they found not. Then he said to them: O foolish, and slow of heart to believe in all things which the prophets have spoken.

Ought not Christ to have suffered these things, and so to enter into his glory? And beginning at Moses and all the prophets, he expounded to them in all the scriptures, the things that were concerning him. And they drew nigh to the town, whither they were going: and he made as though he would go farther. But they constrained him; saying: Stay with us, because it is towards evening, and the day is now far spent. And he went in with them. And it came to pass, whilst he was at table with them, he took bread, and blessed, and brake, and gave to them.

And their eyes were opened, and they knew him: and he vanished out of their sight. And they said one to the other: Was not our heart burning within us, whilst he spoke in this way, and opened to us the scriptures?
\end{quotationx}


It is interesting to note that Luke 24, 13-32 is voluntarily doing a parallel with Genesis 18, 1-33. The structure is the same\footnote{Nota Bene: Thanks are given to the research of Jean-Jacques Lavoie, who formulated the basis of the meditation on the subject.
\begin{enumerate}
\item Arrival of the characters
\item Hospitality towards them
\item Exchange of thoughts
\item Context of feast
\item Superior knowledge given by host
\item Sudden disappearance of characters
\end{enumerate}}. Also, there is a lot of words which are used that parallel the one in both texts («ôphtè», «ophtalmoi», «kai idou», «euriskein», «enantion», «lambanein arton», «kata-klinein», etc.). The reason is to emphasis both tradition and the fact that the death and resurrection of Christ is a good news («evangelion»), just as with Sarah and Abraham, and a primordial return to creation.

The first thing that sets the mood is the double irony presented to the reader (or hearer): the two men don't recognize the truth from both the women and Jesus, just in front of them. The women, even if they represent in front of the law less than man (you need two women for one man in justice), believed the angels (Luke 24, 5-12) even if they didn't saw Jesus. They went beyond reason. Those man cannot open their eyes, open their heart (seat of the intellect even here, see Luke 24, 25; 32) to higher truths, with their «faces downcast». Their earthly feelings cloud them from seeing the truth.

The second thing that is important with the teaching of Luke is the kerygmatic doctrine hidden in the story. It is twofold. First of all, the meaning is that of the traditional liturgy of the Christian Church: it is Sunday, it uses mediation of the community, the Exalted One has the initiative and culminates in the initatic ritual of the Eucharist; it also uses the same structure: teachings from the Scriptures (old and new), then an interpretation and actualization and finally the Eucharist. Second, it is a summary of the teachings of Jesus Christ in the Gospel of Luke, even paralleling the structure of Luke 10 (the Samaritan). The Eucharist uses words from the separation of bread for the Multiplicity earlier in Luke: it marks the Christ as the sustained both of the flesh (Luke 9, 10-17) and the spirit (Luke 24, as the «opener of the scriptures»).

But thirdly and most importantly, this story asks the question: how can we move from non-faith to faith? What is union with Christ? Both disciples gave up, closed their eyes to the truth, spoke only of what they knew on the rational level. They knew a lot of facts; they had proofs (Scriptures, women, etc.); yet they were still closed to truth and faith. Their hearts (because of burning passions), seat of the intellect, was kept hidden from them. But the Logos, both as a scriptural body (teachings, liturgy) and Eucharistic body (rites, initiation) opened their hearts: it was an obligated passage. They thought Jesus as a prophet at first, powerful in deeds and speech; then he was a Messiah coming to save the Kingdom of God on earth; then he was revealed as what He is. The only way they could see him is to go through the process, to walk the Way\footnote{There is a lot of vocabulary in Luke which alludes to the road, the way, or walking: «eisodo», «odos» and «exodos» mainly. The beginning is marked by the «eisodo», and his death by «exodos». Even the following of Christ is described as a way, a path.} with Him. They had to be initiated. Then, and only then, their eyes were opened to His glory, through the sacramental rite, which would now have to be repeated each Sunday; but also at each meal, at each moment of their lives to live in Him, to be dead to the world and alive in Him, just as He did (see Saint Paul letter to the Romans 6, 3-11). Why is it that at that precise moment he disappeared? Because His mode of presence is absence. God reveals Himself in absence, for all times and ages; it is still his mode in our time\footnote{For those who wait for a physical sign from God, a presence in the historical dimension, etc. Luke made himself quite clear: God hasn't changed his mode of presence; we are not less or more suited to Union with Him than the first christians were. The rite of the Eucharist is the mode of presence of God in the world.}. Because He is beyond-being, beyond-existence. He is, yet he is not, because as the fathers have once said: « If God exist, we do not exist. If we exist, God do not exist. » His presence is absence; His absence is Presence.


\flrightit{Posted on 2013-03-30 by David}
